\documentclass[12pt]{article}
\usepackage[utf8]{inputenc}
\usepackage[T1]{fontenc}
\usepackage{hyperref}
\title{Tembo Language Resources for K-12 Curriculum Development}
\author{Compiled by Claude Sonnet 5 (claude-sonnet-5) for Andr\'as Kornai}
\date{}
\begin{document}
\maketitle

\section{Introduction}
Tembo (Motembo, Litembo; ISO 639-3 tmv; Glottocode temb1272) is a Bantu language of the Democratic Republic of Congo. Glottolog places it in the Motembo-Kunda group of the Ngiri branch. Sources give about 5{,}000 speakers (Global Recordings Network) or 18{,}000 people in the Tembo-Equateur group (Joshua Project). It must not be confused with Kitembo (ISO 639-3 tbt) of the Kivu provinces, which has a Bible translation; those resources are not listed here.

\section{Primary Educational Resources}
No primers, textbooks or school programmes for tmv were found in search results.

\section{Linguistic and Reference Materials}
\begin{itemize}
\item Glottolog entry for Tembo: \url{https://glottolog.org/resource/languoid/id/temb1272.bigmap.html}
\item Wikipedia, Budza language, which discusses Tembo and its neighbours: \url{https://en.wikipedia.org/wiki/Budza_language}
\end{itemize}

\section{Religious and Community Resources}
\begin{itemize}
\item Joshua Project language page for tmv: \url{https://joshuaproject.net/index.php/languages/tmv}
\end{itemize}

\section{Digital and Audio Resources}
\begin{itemize}
\item Global Recordings Network page for tmv, which may hold audio. The search did not show its contents: \url{https://globalrecordings.net/en/language/tmv}
\end{itemize}

\section{Curriculum Development Context}
Almost nothing is documented for this language. Any curriculum would start from basic documentation and orthography work, and would probably build on neighbouring Ngiri-branch languages and on the national languages used in schools.
\end{document}